\documentclass[11pt,a4paper]{article}

% ── Codificación y fuentes ────────────────────────────────────────────────────
\usepackage[utf8]{inputenc}
\usepackage[T1]{fontenc}
\usepackage{lmodern}
\usepackage[spanish,es-noshorthands]{babel}

% ── Geometría y espaciado ─────────────────────────────────────────────────────
\usepackage[top=2.5cm, bottom=2.5cm, left=2.8cm, right=2.8cm]{geometry}
\usepackage{setspace}
\setstretch{1.10}
\usepackage{parskip}

% ── Matemáticas ───────────────────────────────────────────────────────────────
\usepackage{amsmath}
\usepackage{amssymb}
\usepackage{siunitx}

% ── Circuitos ────────────────────────────────────────────────────────────────
\usepackage[european, straightvoltages]{circuitikz}

% ── Tablas ────────────────────────────────────────────────────────────────────
\usepackage{booktabs}
\usepackage{array}
\usepackage{enumitem}

% ── Colores y cajas ───────────────────────────────────────────────────────────
\usepackage[dvipsnames]{xcolor}
\usepackage{tcolorbox}
\tcbuselibrary{skins, breakable}

% ── Cabeceras y pies ──────────────────────────────────────────────────────────
\usepackage{fancyhdr}
\usepackage{lastpage}
\pagestyle{fancy}
\fancyhf{}

\fancyhead[L]{\small\textbf{ELECTRÓNICA} --- Examen}
\fancyhead[R]{\small\itshape Examen de Electrónica: Leyes de Kirchhoff}
\fancyfoot[C]{\small Página \thepage\ de \pageref{LastPage}}
\renewcommand{\headrulewidth}{0.4pt}

% ── Hiperenlaces ──────────────────────────────────────────────────────────────
\usepackage[colorlinks=true, linkcolor=NavyBlue, urlcolor=NavyBlue]{hyperref}

% ── Colores personalizados ────────────────────────────────────────────────────
\definecolor{azulTitulo}{HTML}{1B2A6B}
\definecolor{grisFondo}{HTML}{F5F5F5}
\definecolor{verdePuntaje}{HTML}{1A6B2F}

% ── Estilos de cajas ─────────────────────────────────────────────────────────
\tcbset{
  cajaInstrucciones/.style={
    enhanced, breakable,
    colback=grisFondo, colframe=gray!50,
    fonttitle=\bfseries, coltitle=black,
    top=6pt, bottom=6pt, left=8pt, right=8pt,
  },
}

% ──────────────────────────────────────────────────────────────────────────────
\begin{document}

% ══ PORTADA / ENCABEZADO ═══════════════════════════════════════════════════════
\begin{center}
  {\Large\bfseries\color{azulTitulo} ELECTRÓNICA}\\[4pt]
  {\large Examen de Razonamiento y Cálculo}\\[8pt]
  \rule{\linewidth}{1.2pt}\\[6pt]
  {\LARGE\bfseries Examen de Electrónica: Leyes de Kirchhoff}\\[6pt]
  \rule{\linewidth}{0.6pt}\\[6pt]
  {\large \textbf{Dificultad:} Intermedia \hfill
   \textbf{Duración:} 90 minutos} \\[4pt]
  {\large \textbf{Fecha:} \rule{4cm}{0.2pt} \hfill
   \textbf{Estudiante:} \rule{6cm}{0.2pt}} \\[6pt]
  \rule{\linewidth}{0.2pt}
\end{center}

% ══ DATOS DEL EXAMEN ═══════════════════════════════════════════════════════════
\vspace{4pt}
\begin{tcolorbox}[cajaInstrucciones, title={Instrucciones}]
Resolver cada pregunta paso a paso, justificando los cálculos y utilizando la notación y fórmulas correctas. Se permite el uso de una calculadora científica y un formulario personal de una hoja.
\medskip
\textbf{Material permitido:} Calculadora científica, formulario personal de una hoja.
\end{tcolorbox}

% ══ INDICACIÓN DE PUNTAJE ═════════════════════════════════════════════════════
\vspace{4pt}
\begin{center}
  \begin{tcolorbox}[
    enhanced,
    colback=white, colframe=azulTitulo,
    width=0.6\linewidth,
    boxrule=1.5pt,
    halign=center, valign=center,
    top=4pt, bottom=4pt,
  ]
    \centering
    \textbf{Puntaje total: 10.0 puntos} \\
    \small\textit{Responda cada pregunta en el espacio provisto. \\
    Debe mostrar todo el desarrollo matemático para obtener puntaje completo.}
  \end{tcolorbox}
\end{center}

% ══ PREGUNTAS ══════════════════════════════════════════════════════════════════
\section*{Preguntas}


% ── Pregunta 1 ──────────────────────────────────────────────────────
\subsection*{Pregunta 1 \normalfont\normalsize\textit{(2.0 puntos)}}
En el circuito mostrado, determine la corriente $I_1$ que atraviesa la resistencia $R_1$ de $2Omega$ y la corriente $I_2$ que atraviesa la resistencia $R_2$ de $3Omega$, sabiendo que la fuente de voltaje es de $12V$. Utilice la Primera Ley de Kirchhoff para encontrar las corrientes desconocidas. $R_3 = 4Omega$.

  \begin{tcolorbox}[colback=gray!5, colframe=gray!40, fonttitle=\bfseries, title={Diagrama sugerido}]
    Un circuito con una fuente de voltaje de $12V$, resistencias $R_1 = 2Omega$, $R_2 = 3Omega$, y $R_3 = 4Omega$ conectadas en serie.
  \end{tcolorbox}

\vspace{2mm}
\rule{\linewidth}{0.2pt}
\vspace{3cm}  % espacio para respuesta



% ── Pregunta 2 ──────────────────────────────────────────────────────
\subsection*{Pregunta 2 \normalfont\normalsize\textit{(2.2 puntos)}}
Un circuito tiene dos fuentes de voltaje en serie, $V_1 = 10V$ y $V_2 = 6V$, con resistencias $R_1 = 2Omega$ y $R_2 = 3Omega$ también en serie. Determine la corriente $I$ que circula por el circuito y el voltaje $V_{out}$ en los bornes de $R_2$ utilizando las Leyes de Kirchhoff.

  \begin{tcolorbox}[colback=gray!5, colframe=gray!40, fonttitle=\bfseries, title={Diagrama sugerido}]
    Un circuito con dos fuentes de voltaje $V_1$ y $V_2$ en serie con resistencias $R_1$ y $R_2$ también en serie.
  \end{tcolorbox}

\vspace{2mm}
\rule{\linewidth}{0.2pt}
\vspace{3cm}  % espacio para respuesta



% ── Pregunta 3 ──────────────────────────────────────────────────────
\subsection*{Pregunta 3 \normalfont\normalsize\textit{(2.5 puntos)}}
En el circuito mostrado, se desconocen las corrientes $I_1$, $I_2$, e $I_3$. Las resistencias son $R_1 = 4Omega$, $R_2 = 6Omega$, y $R_3 = 3Omega$. Las fuentes de voltaje son $V_1 = 12V$ y $V_2 = 8V$. Utilice las Leyes de Kirchhoff para determinar estas corrientes.

  \begin{tcolorbox}[colback=gray!5, colframe=gray!40, fonttitle=\bfseries, title={Diagrama sugerido}]
    Un circuito con dos fuentes de voltaje $V_1$ y $V_2$, y tres resistencias $R_1$, $R_2$, $R_3$ configuradas en dos lazos con un nodo común.
  \end{tcolorbox}

\vspace{2mm}
\rule{\linewidth}{0.2pt}
\vspace{3cm}  % espacio para respuesta



% ── Pregunta 4 ──────────────────────────────────────────────────────
\subsection*{Pregunta 4 \normalfont\normalsize\textit{(2.0 puntos)}}
Un circuito tiene dos resistencias $R_1 = 10Omega$ y $R_2 = 15Omega$ conectadas en paralelo, con una fuente de voltaje $V = 20V$ conectada a ellas. Determine la corriente total $I_{total}$ que sale de la fuente y las corrientes $I_1$ e $I_2$ que circulan por cada resistencia.

  \begin{tcolorbox}[colback=gray!5, colframe=gray!40, fonttitle=\bfseries, title={Diagrama sugerido}]
    Un circuito con una fuente de voltaje $V$ y dos resistencias $R_1$ y $R_2$ en paralelo.
  \end{tcolorbox}

\vspace{2mm}
\rule{\linewidth}{0.2pt}
\vspace{3cm}  % espacio para respuesta



% ── Pregunta 5 ──────────────────────────────────────────────────────
\subsection*{Pregunta 5 \normalfont\normalsize\textit{(3.3 puntos)}}
En un circuito con varias ramas y fuentes de voltaje, se aplican las Leyes de Kirchhoff para encontrar las corrientes desconocidas. Considerando el circuito con fuentes de voltaje $V_1 = 15V$ y $V_2 = 10V$, y resistencias $R_1 = 5Omega$, $R_2 = 8Omega$, $R_3 = 12Omega$, determine las corrientes $I_1$, $I_2$, e $I_3$.

  \begin{tcolorbox}[colback=gray!5, colframe=gray!40, fonttitle=\bfseries, title={Diagrama sugerido}]
    Un circuito con dos fuentes de voltaje $V_1$ y $V_2$, y tres resistencias $R_1$, $R_2$, $R_3$ configuradas en dos lazos con un nodo común.
  \end{tcolorbox}

\vspace{2mm}
\rule{\linewidth}{0.2pt}
\vspace{3cm}  % espacio para respuesta



% ══ FIN DEL EXAMEN ═════════════════════════════════════════════════════════════
\vspace{12pt}
\begin{center}
  \rule{0.6\linewidth}{0.4pt} \\[6pt]
  \textbf{--- Fin del examen ---}
\end{center}

\vfill
\begin{center}
  \footnotesize\color{gray}
  Examen generado automáticamente por el sistema de inteligencia artificial (llama-3.3-70b-versatile) \\
  ELECTRÓNICA --- Sistema de Evaluación Académica \\
  Generado el 2026-06-28
\end{center}

\end{document}
